\subsection*{38. The Classical Evidence Was Influential; and Later Corrected}

Frank Schmidt and John Hunter's 1998 review synthesized decades of personnel-selection research. It popularized strong estimates for general mental ability, work samples, and structured interviews, including high estimated validity for some combinations [R1]. These findings helped shift attention from unstructured judgment to systematic assessment. Many practitioner tables still reproduce their numbers, often without explaining the statistical adjustments behind them.

Paul Sackett, Charlene Zhang, Christopher Berry, and Filip Lievens revisited these methods in 2022, arguing that commonly used corrections for range restriction had systematically inflated many estimates. Their revised analyses generally retained the usefulness of selection procedures but lowered estimated relationships, often by approximately 0.10 to 0.20 correlation points, with \textbf{structured interviews emerging as the highest-ranked procedure} in their reassessment [R2]. Follow-up work examined how these changes affect selection-system design [R5].

The appropriate lesson is neither to discard cognitive assessment nor to treat interviews as infallible. It is that \textbf{measurement research evolves}, the correction method matters, and an entrepreneur should prefer several defensible sources of job-relevant evidence over an overconfident ranking of instruments by a single historical coefficient.

\subsection*{39. A Validity Ranking Cannot Tell You Which Test to Use Without the Job}

A general ranking hides important moderators. A work sample may be excellent for a role requiring immediately demonstrable competence and less appropriate for an entry-level trainee expected to learn unfamiliar tasks. A verbal assessment may matter in a negotiation role but add little if it measures accent rather than reasoning. A structured interview can evaluate conflict resolution, teaching judgment, or commercial ethics, but only if the prompts and scoring correspond to the job.

The right evidence bundle depends on the \textbf{selection problem}: number of candidates, required competencies, training available, job stakes, candidate burden, and current evidence. An organization with 5 shortlisted candidates can use detailed simulations. An organization with 5,000 applicants needs an accessible, standardized initial screen before labor-intensive evaluations.

Research supports building a structured system. It does not authorize selecting a method from a generic league table and ignoring the role. A new educational venture should ask which faculty or reviewer behaviors most strongly affect its particular promise and then choose assessments that observe them.

\subsection*{40. Range Restriction Distorts the Evidence We See}

Suppose a company hires only applicants with cognitive scores above a demanding cutoff. Later, it observes a weak relationship between scores and performance among those hired. This does not necessarily mean cognitive ability has no predictive value across the original applicant population. The company has compressed the range of observed scores. On the other hand, statistical corrections for this compression can be too aggressive if their assumptions are wrong, as the modern reanalysis emphasizes.

The practical complication is that \textbf{rejected candidates have no on-the-job performance data}. A founder cannot directly know whether an excellent rejected candidate would have succeeded. This creates selection bias when validating a hiring process from its own employees. Promotion decisions face similar problems: only people given opportunities can demonstrate performance in those opportunities.

The organization should record assessment distributions, selection decisions, and eventual outcomes with appropriate privacy safeguards. Interpret observed relationships cautiously and use professional expertise for formal validation. Do not ``correct'' correlations with a spreadsheet formula merely because a textbook mentions range restriction.

\subsection*{41. Measurement Error Creates Both Individual and Organizational Uncertainty}

Suppose a candidate receives an observed assessment score $X = T + e$, where $T$ represents a stable capability component and $e$ represents measurement error or temporary variation. Under a classical test-theory model, reliability can be expressed conceptually as the proportion of observed-score variance attributable to the true-score component, given the model's assumptions:
\begin{equation}
\rho_{XX'} = \frac{\operatorname{Var}(T)}{\operatorname{Var}(X)}.
\end{equation}

This mathematical framing is useful because 2 candidates with scores of 81 and 83 may be indistinguishable at the precision of the test. Treating the 2-point gap as decisive can amount to selecting based on noise. Similarly, a candidate who receives 1 poor interview score may deserve more evidence when the measured competency is central and the result conflicts sharply with other observations.

Score precision should influence how a decision is made. Use ranges, evidence summaries, and discussion of uncertainty rather than arbitrary decimal places. The goal is not to eliminate ambiguity by formatting it beautifully.

\subsection*{42. Independent Evidence Is More Useful Than Repeated Impressions}

3 interviewers can each report that a candidate is ``confident,'' but if they all observed the same presentation style, the employer has not gathered 3 independent kinds of evidence. This is correlated measurement. More assessors do not automatically improve prediction if they repeat the same questions, hear the same reputation cues, or influence 1 another before scoring.

A stronger portfolio measures different role-relevant attributes through different mechanisms: a realistic task for applied competence; a structured behavioral interview for handling difficult situations; a reference for historical reliability; and perhaps a well-designed reasoning assessment where learning and problem solving are essential. The instruments will still correlate, but their \textbf{incremental validity} can be more meaningful than another general conversation.

The founder should build an evidence map showing which important competencies are assessed by which methods. Overlap is appropriate for high-risk requirements, especially when 1 method is susceptible to rehearsal. But unnecessary repetition adds burden without certainty. \textbf{More data are not always more information.}
